\section{Introduction}
\label{sec:introduction}

\subsection{Background and Motivation}
\label{sec:background_and_motivation}

\IEEEPARstart{R}{eaching} a net-zero economy by 2050 requires extensive electrification of the energy system. During the past two decades, renewable generation, particularly wind and photovoltaic power, has grown steadily and is expected to accelerate towards 2050~\cite{international_energy_agency_world_2023}. At the same time, decarbonization targets require the electrification of transport, heating, and industrial processes, replacing fossil fuels with electricity. In Europe, studies project that the total electricity demand will triple by 2050~\cite{energyville_paths_2023}. Consequently, future power systems will be largely dominated by power-electronic interfaced resources on both the supply and demand sides. This transition fundamentally changes the harmonic behavior of power systems. Traditional grids, dominated by synchronous machines and mostly linear loads, exhibited limited harmonic distortion. Modern systems, however, increasingly rely on converters whose nonlinear characteristics generate a broad spectrum of harmonic currents. The aggregation of numerous distributed converters creates new challenges in maintaining harmonic voltage quality across transmission and distribution networks.

\begin{table*}[!t]
    \centering
    \caption{Literature review on harmonic hosting capacity, network harmonic impedance uncertainty, bounded-uncertainty harmonic analysis and robust optimization over network parameters, in order of discussion, enumerating the uncertain features~(a)--(e). The entries marked $\star$ are prior work of the present authors.}
    \label{tab:literature_review}
    \begin{tabular}{ c | c | c | c c | c | c c c c c | r r }
        \multirow{2}{*}{reference} & \multirow{2}{*}{objective$^{a}$} & \multirow{2}{*}{method$^{b}$} & \multicolumn{2}{c |}{network} & \multirow{2}{*}{uncertainty$^{c}$} & \multicolumn{5}{c |}{uncertain features} & \multicolumn{2}{c}{tractability$^{d}$} \\
                                            &       &       & model & phases &       & (a)  & (b)   & (c)   & (d)   & (e)   & nodes & minutes \\ \hline
        \cite{noauthor_iec61000-3-6_2008}   & hhc   & ae    & nodal & bal.  & an.   &\cmark &\xmark &\xmark &\xmark &\cmark & --    & --    \\
        \cite{santos_considerations_2015}   & hhc   & ae    & nodal & bal.  & an.   &\cmark &\cmark &\xmark &\xmark &\xmark & --    & --    \\
        \cite{sun_distribution_2022}        & hhc   & cc    & full  & bal.  & scen. &\xmark &\cmark &\xmark &\cmark &\cmark & 16    & 30    \\
        \cite{nakhodchi_impact_2023}        & hhc   & ae    & nodal & bal.  & mc    &\cmark &\xmark &\xmark &\xmark &\xmark & 64    & --    \\
        \cite{ismael_probabilistic_2019}    & hhc   & mh    & full  & bal.  & mc    &\cmark &\cmark &\cmark &\cmark &\cmark & --    & --    \\
        \cite{emini_study_2026}             & hhc   & pf    & full  & bal.  & mc    &\xmark &\xmark &\xmark &\cmark &\xmark & --    & --    \\ \arrayrulecolor{gray!20} \hline
        \cite{testa_probabilistic_2002}     & hpf   & pf    & full  & bal.  & mc    &\cmark &\xmark &\cmark &\cmark &\xmark & 13    & --    \\
        \cite{li_probabilistic_2018}        & hra   & rma   & full  & bal.  & mc    &\xmark &\xmark &\cmark &\cmark &\xmark & 14    & 4.4   \\
        \cite{li2020}                       & hpf   & pf    & full  & bal.  & mc    &\cmark &\xmark &\cmark &\xmark &\xmark & 14    & --    \\
        \cite{barakou_parameter_2023}       & hra   & fs    & full  & --    & scen. &\xmark &\xmark &\cmark &\cmark &\xmark & --    & --    \\ \arrayrulecolor{gray!20} \hline
        \cite{romero_frequency_2018}        & hra   & fs    & full  & bal.  & pos.  &\xmark &\xmark &\xmark &\cmark &\xmark & 14    & --    \\
        \cite{huang_newton_2017}            & hpf   & pf    & full  & bal.  & int.  &\cmark &\xmark &\xmark &\xmark &\xmark & 30    & --    \\
        \cite{desousa_interval_2024}        & hpf   & pf    & full  & bal.  & int.  &\cmark &\xmark &\xmark &\cmark &\xmark & 33    & $<$\,0.01 \\
        \cite{lin_extended_2017}            & hra   & ae    & full  & bal.  & rob.  &\xmark &\xmark &\cmark &\xmark &\xmark & 81    & $<$\,0.001 \\ \arrayrulecolor{gray!20} \hline
        \cite{christakou_voltage_2018}      & vc    & ro    & lin.  & bal.  & rob.  &\xmark &\xmark &\cmark &\xmark &\xmark & 34    & --    \\
        \cite{liu_feasible_2023}            & fr    & ro    & lin.  & unb.  & rob.  &\xmark &\xmark &\cmark &\xmark &\xmark & --    & 8     \\ \arrayrulecolor{black} \hline
        $\star$~\cite{vanacker_hosting_2026}& hhc   & soc   & full  & bal.  & --    &\xmark &\xmark &\xmark &\xmark &\xmark & 1880  & $<$\,6 \\
        $\star$~\cite{vanacker_shhc}        & hhc   & qcqp/soc & full & bal. & pce  &\cmark &\xmark &\xmark &\xmark &\cmark & 30    & $<$\,0.04 \\ \arrayrulecolor{black} \hline
        \multirow{2}{*}{This work}          & \multirow{2}{*}{hhc}   & \multirow{2}{*}{soc}   & \multirow{2}{*}{full}  & \multirow{2}{*}{bal.}  & an./rob. &\cmark &\xmark &\cmark &\xmark &\cmark & 30$^{*}$ & $<$\,3 \\ 
        &&&&& rob. & \xmark & \xmark & \cmark &\xmark & \xmark & 30$^{**}$ & $<$\,8 \\
    \end{tabular}

    \vspace{2pt}
    \parbox{\textwidth}{\scriptsize
    $^{a}$Objective: (hhc)~harmonic hosting capacity, (hpf)~harmonic power flow, (hra)~harmonic resonance or impedance analysis, (vc)~voltage control and (fr)~feasible region of distributed energy resources, both at the fundamental frequency.
    $^{b}$Mathematical methods: (ae)~analytical expression, (pf)~power flow, (fs)~frequency scan, (rma)~resonance mode analysis, (mh)~metaheuristic optimization, (cc)~chance-constrained optimal power flow, (ro)~robust counterpart, (soc)~second-order cone optimal power flow, (qcqp)~quadratically constrained quadratic program. Network model: (nodal)~equivalent at one node, (full)~full network, (lin.)~linearized full network. Phases: (bal.)~balanced, (unb.)~unbalanced.
    $^{c}$Treatment of uncertainty: (mc)~Monte Carlo simulation, (an.)~analytical or moment-based, (scen.)~scenario-based, (pos.)~possibility theory, (int.)~interval or affine arithmetic, (pce)~polynomial chaos expansion, (rob.)~robust, i.e., guaranteed over a set.
    $^{d}$Tractability expressed as the number of nodes and the solve time in minutes. $^{*}$Positive- and negative-sequence harmonics 5--25, 23 uncertain physical inputs, on 10 threads. $^{**}$Zero-sequence variants without recombination, harmonics 2--50, on 16 threads.}
\end{table*}

A key challenge arises from the stochastic nature of the harmonic impedance of the network. System operators describe cables and overhead lines by their sequence resistance, reactance and susceptance at the fundamental frequency, taken from datasheets or commissioning measurements, and harmonic studies extrapolate these values to higher frequencies with scaling rules. In reality, the harmonic impedance follows from the physical inputs of each cable and line, e.g., the resistivities and temperatures of the conductors and screens, the permittivity of the insulation, the dimensions within their manufacturing tolerances, the laying geometry and the earth. ``In general, most of the physical parameters of a cable such as the resistivity, the permittivity and $\tan{\delta}$ are frequency- and temperature-dependent, and such data are not even available from a cable manufacturer. The same is true for the earth resistivity and permittivity.''~\cite[\S\,2.5.2.6]{ametani_cable_2015}. A cable project adds its own variability: the conductor temperature follows the load, the earth resistivity varies by orders of magnitude along the route and over the seasons, and the installed cables depart from their design within the acceptance limits of the standards. The harmonic impedances near a parallel resonance are the most sensitive to these inputs, and the scaling rules miss their joint effect. As the penetration of the converter increases, system operators need transparent metrics to assess whether additional units can be connected without exceeding harmonic limits. Harmonic hosting capacity provides such a metric by defining the maximum allowable harmonic emissions while ensuring compliance with power-quality standards.  The harmonic hosting capacity must therefore hold for every network that the physical inputs allow. However, most existing methods rely on impedance data at the fundamental frequency, extrapolated towards higher harmonics with scaling rules, leading to overly conservative or insufficiently protective results.

\subsection{Literature Review}
\label{sec:literature_review}

A growing body of literature addresses the harmonic hosting capacity of a power system, the uncertainty of the network harmonic impedance, harmonic analysis under bounded rather than probabilistic uncertainty, and optimization with a guarantee over uncertain data. Table~\ref{tab:literature_review} summarizes the publications discussed below, characterized by their objective, mathematical method, network representation, treatment of uncertainty and tractability, and by which of the following features they treat as uncertain:
\begin{enumerate*}
    \item [(a)] unit harmonic injection,
    \item [(b)] background harmonic distortion,
    \item [(c)] network parameters,
    \item [(d)] network exploitation, and
    \item [(e)] harmonic limits.
\end{enumerate*}

\subsubsection{Harmonic hosting capacity}

The standard IEC~61000-3-6 allocates planning levels over the installations connected to a point of common coupling and aggregates their emissions with an indicative summation law~\cite{noauthor_iec61000-3-6_2008}. Santos et al.\ used the summation law to determine the harmonic hosting capacity behind a Norton equivalent with background distortion of unknown phase~\cite{santos_considerations_2015}. Sun and Harrison proposed the first optimization-based stochastic method, with percentile-based harmonic distortion limits as chance constraints within an ac optimal power flow~\cite{sun_distribution_2022}. Nakhodchi et al.\ showed for a fast charging station for electric vehicles that the estimated hosting capacity depends strongly on how the medium-voltage network and the remote loads are modeled~\cite{nakhodchi_impact_2023}. Ismael et al.\ designed a filter that maximizes the harmonic-constrained hosting capacity of a photovoltaic system, with a metaheuristic around a Monte Carlo simulation over, among others, the filter parameters~\cite{ismael_probabilistic_2019}. At transmission level, a Monte Carlo study over the system loading found the harmonic hosting capacity to be governed by resonance-sensitive impedance~\cite{emini_study_2026}. The present authors determined an exact lower bound for the equitable harmonic hosting capacity with a second-order cone formulation of the full network~\cite{vanacker_hosting_2026}, showed that underground cables change it substantially~\cite{vanacker_impact_2025}, and extended it to stochastic harmonic current angles with chance constraints~\cite{vanacker_shhc}. Except for the filter in~\cite{ismael_probabilistic_2019} all of these methods take the network parameters as known.

\subsubsection{Uncertainty of the network harmonic impedance}
The joint CIGRE and CIRED guide stresses that the harmonic impedance changes continuously with load and network configuration and limits the validity of network models to about the 20th harmonic~\cite{robert_guide_1997}. Testa et al.\ built a probabilistic model of the nodal impedances from a 5\,\% tolerance on the inductances and capacitances, and concluded that harmonic voltage assessments that ignore the impedance variability give unreliable results~\cite{testa_probabilistic_2002}. Li et al.\ propagated a 10\,\% variation of the line, transformer and capacitor parameters through resonance mode analysis~\cite{xu_harmonic_2005,huang_application_2007} by Monte Carlo simulation, and obtained probabilistic resonance frequency bands and modal impedances~\cite{li_probabilistic_2018}. The same group separated aleatory from epistemic sources~\cite{li2020}. For offshore wind, Barakou et al.\ found that network and plant contingencies dominate the harmonic amplification, that the choice of cable model is noticeable and that a manufacturing tolerance of 8\,\% on the cable capacitance has a minor effect~\cite{barakou_parameter_2023}. Babaev et al.\ estimated the cable harmonic impedance online with an uncertainty budget, as its theoretical derivation neglects the installation~\cite{babaev_online_2020}. These works sample or scan the network parameters, but none of them optimizes the hosting capacity over their uncertainty.

\subsubsection{Harmonic analysis under bounded uncertainty}

Where probability distributions are not available, uncertainty is described by bounds. Romero et al.\ introduced the harmonic load flow based on possibility theory~\cite{romero_harmonic_2011} and extended the frequency scan with it to the possibility of a parallel resonance~\cite{romero_frequency_2018}. Interval and affine arithmetic bound the power flow solution for interval data~\cite{vaccaro_affine_2010} and were carried to harmonics. Huang et al.\ proposed a Newton-Raphson interval harmonic power flow in affine form~\cite{huang_newton_2017} and de Sousa et al.\ an interval harmonic power flow for microgrids based on a second-order Taylor expansion~\cite{desousa_interval_2024}. In all of these, the bounded quantities are the injections, while the network admittance is fixed. The paper presented by Lin et al.\ is the exception, based on the extended Gersgorin theorem, they derived a parameter domain that is free of harmonic resonance, which is conservative by construction~\cite{lin_extended_2017}. These methods bound a response or a feasible domain for given injections, but do not determine the hosting capacity.

\subsubsection{Robust optimization}

Robust optimization requires a constraint to hold for every realization of uncertain data within a set, and replaces the uncertain program by its robust counterpart~\cite{ben-tal_robust_1999}. Where no tractable counterpart exists, cutting-set methods alternate between a master problem over finite scenario sets and a worst-case analysis~\cite{mutapcic_cutting-set_2009}. In power systems, robust optimization has mainly addressed the uncertainty of renewable generation and load~\cite{roald_power_2023}, e.g., in robust ac and chance-constrained optimal power flow and hosting capacity~\cite{louca_robust_2019,lee_robust_2021,lubin_robust_2016,wang_distributed_2016}. Few works place the network itself in the uncertainty set. Christakou et al.\ controlled the voltage of a distribution network with line resistances that are bounded but of unknown distribution, due to conductor temperature and incomplete asset data, and linearized the network with sensitivity coefficients~\cite{christakou_voltage_2018}. Liu and Ma determined the robust feasible region of distributed energy resources for interval line impedances in a linear unbalanced power flow~\cite{liu_feasible_2023}. All of these are linear or linearized models at the fundamental frequency, whereas the network parameters enter the harmonic hosting capacity through the inverse of the harmonic admittance matrix, which is far from concave near a parallel resonance.

\subsubsection{Research gaps}

No existing method simultaneously achieves:
\begin{enumerate*}
    \item [(i)] an optimization of the harmonic hosting capacity with a guarantee over the uncertainty of the network parameters, i.e., robustly rather than by sampling,
    \item [(ii)] network parameters derived from the physical construction of the cables and lines and their frequency dependence, rather than from scaled fundamental-frequency data, and
    \item [(iii)] the joint treatment of both uncertainties, the angles in probability and the network in the worst case.
\end{enumerate*}
Moreover, the zero-sequence harmonics, for which the transformer windings and earthing set the network, are absent from the optimization-based methods.

% Table II (Section II) is defined here so that, as a double-column float, it lands at the bottom of the page on which Section II starts
\begin{table*}[!b]
    \centering
    \caption{Cables, overhead lines and earth of the IEEE 30-bus case: layout, and nominal value and interval of the 11 most influential of the 23 uncertain inputs, with the largest relative change of a sequence parameter at \textcolor{black}{50, 650 or 1250\,Hz} when the input alone moves to an end of its interval; $z_{0}$, $z_{1}$ and $y_{0}$, $y_{1}$ are the zero- and positive-sequence series impedance and shunt admittance per unit length. }
    \label{tab:parameter_box}
    \footnotesize\setlength{\tabcolsep}{4pt}
    \begin{tabular}{ >{\raggedright\arraybackslash}p{0.42\textwidth} >{\raggedright\arraybackslash}p{0.20\textwidth} l r }
        quantity & nominal & interval & largest change \\ \hline
        \multicolumn{4}{l}{\textit{cables, 33\,kV, 22 links in touching trefoil, center 0.8\,m deep, screens bonded at both ends}} \\
        aluminum conductor cross section~\cite{iec60228_2023} & 400--800\,mm$^{2}$ & by link & -- \\
        conductor DC resistance at 20\,$^{\circ}$C~\cite{iec60228_2023} & 0.0367--0.0778\,$\Omega$/km & $-8$\,\% to nominal & 4.4\,\% $z_{1}$ \\
        conductor diameter~\cite{iec60228_2023} & 23.0--34.2\,mm & nominal to $+0.4$\,mm & 1.3\,\% $y_{0}$ \\
        XLPE insulation thickness~\cite{iec60502_2_2014} & 8.0\,mm & 7.6 to 8.4\,mm & 4.4\,\% $y_{0}$ \\
        XLPE relative permittivity~\cite{iec60287_1_1_2006} & 2.4 & 2.1 to 2.5 & 12.5\,\% $y_{1}$ \\
        screen DC resistance at 20\,$^{\circ}$C, 48 copper wires of 1.35\,mm & 0.263\,$\Omega$/km & 0.251 to 0.263\,$\Omega$/km & 3.8\,\% $z_{0}$ \\
        jacket thickness~\cite{iec60502_2_2014} & 2.85--3.30\,mm & nominal to $+35$\,\% & 2.0\,\% $z_{1}$ \\
        conductor temperature~\cite{iec60502_2_2014} & 20\,$^{\circ}$C & 10 to 90\,$^{\circ}$C & 26\,\% $z_{0}$ \\
        \multicolumn{4}{l}{\textit{overhead lines, 132\,kV, 12 single-circuit lines in delta, span 300\,m}} \\
        ACSR phase conductors, 30 aluminum and 7 steel wires~\cite{en50182_2001} & Lion, Goat, Bear, Panther & by line & -- \\
        phase positions $(x, y)$ at the tower & $(\pm 2.5, 18.0)$, $(0, 22.33)$\,m & -- & -- \\
        earth wire, 7 steel wires of 3.15\,mm~\cite{iec63248_2022}, position & $(0, 29.0)$\,m & -- & -- \\
        reference sag at 20\,$^{\circ}$C, phases and earth wire & 6.0 and 5.4\,m & -- & -- \\
        phase temperature & 20\,$^{\circ}$C & $-5$ to 75\,$^{\circ}$C & 6.9\,\% $z_{1}$ \\
        ambient and earth-wire temperature & 20\,$^{\circ}$C & $-5$ to 45\,$^{\circ}$C & 3.2\,\% $z_{1}$ \\
        earth-wire steel relative permeability & 100 & 50 to 300 & 3.1\,\% $z_{0}$ \\
        \multicolumn{4}{l}{\textit{earth, homogeneous}} \\
        resistivity, logarithmic scale & 100\,$\Omega\cdot$m & 10 to 1000\,$\Omega\cdot$m & 14.6\,\% $z_{0}$ \\
        relative permittivity & 10 & -- & -- \\
        \multicolumn{4}{l}{\textit{12 further inputs}} \\
        semiconducting screens; wire diameters and lay ratios~\cite{en50182_2001,iec62641_2022,iec63248_2022}; metal resistivities; phase steel permeability & & & \textcolor{black}{$\leq 0.65$\,\%}
    \end{tabular}
\end{table*}

\subsection{Contributions and Outline}
\label{sec:contribution_and_outline}

In this paper, a holistic one-shot methodology is presented to determine the steady-state harmonic hosting capacity for an entire phase-balanced three-phase power system, extending~\cite{vanacker_shhc}. The presented method considers uncertainty on both: 
\begin{enumerate}
    \item [(i)] the angle of the injected harmonic current, enabling harmonic recombination, and
    \item [(ii)] cable and overhead-line parameters, 
\end{enumerate}
limiting the resulting individual and total harmonic voltage distortion, as well as the root-mean-square voltages and currents, using chance constraints. The problem is formulated from the perspective of the power system operator, who may anticipate the converter technologies its customers will connect, but not the operating point, parameters and control of the individual units. The angle of the harmonic current is therefore modeled as a random variable per unit. Additionally, the cable and overhead-line parameters follow from physical inputs, e.g., conductor temperatures, manufacturing tolerances, and earth resistances, known to lie within intervals. The harmonic hosting capacity is determined in terms of harmonic current magnitude using a fairness principle, anticipating angle uncertainty, and robust against parameter uncertainty. The robust counterpart over the parameter set is solved by a cutting-set method which gives a robust upper bound on the harmonic hosting capacity. The presented formulation uses the rectangular current-voltage variable space and incorporates:
\begin{itemize}
    \item uncertainty on the angle of the injected harmonic current via the angle model of~\cite{vanacker_shhc}, with exact moments in place of its truncated polynomial chaos expansion,
    \item series impedance and shunt admittance of the cables and overhead lines computed at every harmonic frequency from their physical inputs, with intervals from the acceptance limits of the standards and the operating ranges,
    \item exact representation of the harmonic impedance in the frequency domain and detailed treatment of transformers, in the zero, positive and negative sequence,
    \item enforcement of stochastic harmonic limits on bus voltage and component ampacity through chance constraints, and
    \item equitable distribution of the harmonic hosting capacity following a fairness principle.
\end{itemize}
For fairness criteria that fix the ratios between the allocations of the units, e.g., absolute equality, and a constant risk factor, the robust counterpart is a convex semi-infinite second-order cone program with one pair of variables per harmonic. The models and numerical illustrations are made freely available as part of the \textsc{Julia}-package \textsc{HarmonicPowerModels.jl}\footnote{Available at \url{https://github.com/timmyfaraday/HarmonicPowerModels.jl}.}, with the cable and line parameters from \textsc{LineCableModels.jl}\footnote{Available at \url{https://github.com/Electa-Git/LineCableModelsm.jl}.}~\cite{lcm_2026}.

The remainder of the paper is organized as follows: Section~\ref{sec:parameter_uncertainty} describes the physical inputs of the cables and overhead lines and their intervals. Section~\ref{sec:mathematical_framework} presents the robust counterpart of the stochastic harmonic hosting capacity problem and its solution by a cutting-set method. Section~\ref{sec:numerical_illustration} compares the risk factors and applies the method to the IEEE 30-bus system with physical cable and line inputs, including the zero-sequence harmonics and the transformer earthing. Section~\ref{sec:conclusion} concludes the paper.
